\subsubsection{Reduced-model ML/MAP estimation and restart operations}
\label{sec:reduced-od-estimator}

Phase~10 fits the minimal gravity model through
\path{public_transportation.inference.reduced_od.estimator}.  Maximum
likelihood minimizes the negative log likelihood from
Section~\ref{sec:reduced-od-minimal-gravity}.  MAP uses the same compiled
objective and adds independent Gaussian penalties on the \emph{raw},
unconstrained parameter vector $\bm\eta$:
\[
  Q_{\mathrm{MAP}}(\bm\eta)
  =-\ell(\bm\eta)
   +\frac12\sum_j
    \left(\frac{\eta_j-m_j}{s_j}\right)^2.
\]
The softplus transformation described above is applied after this raw-vector
prior.  Thus a prior intended for a transformed behavioral coefficient must be
converted deliberately rather than being interpreted as a raw Gaussian.
Setting every $s_j=+\infty$ gives an exactly flat penalty and makes MAP
numerically identical to ML.  A finite informative prior generally changes the
estimate, especially in weakly identified directions; neither estimate should
be expected to equal a Bayesian posterior mean.

SciPy L-BFGS-B drives a JIT-compiled JAX value-and-gradient calculation.  The
optimizer dimension is the small number of gravity and production-basis
parameters, not the number of feasible OD cells.  Objective evaluations retain
only grouped demand and compact response arrays.  They cannot call detailed
assignment, reconstruct a full OD vector, or write per-cell iteration output.
Compilation and optimization times, evaluation and iteration counts, final
gradient norm, likelihood, prior contribution, method, and model fingerprint
are recorded separately.

The operational layer in \path{operations.py} defines immutable fit manifests
and atomic JSON checkpoints.  A checkpoint contains the model fingerprint,
method, parameter dimension, latest raw parameters, objective, iteration, and
elapsed time.  Resume first requires exact manifest compatibility; a cache from
a different problem, statistical model, method, or dimension fails closed.
Progress callbacks and checkpoint writes occur only at declared iteration
boundaries.  A deadline produces status \texttt{deadline}, saves a resumable
checkpoint, and never claims successful completion.  Result files likewise
carry an explicit \texttt{complete}, \texttt{deadline}, or \texttt{failed}
status.

On the public 20,000-cell, 2,000-measurement synthetic benchmark, the
two-parameter ML fit compiled in about 0.13 seconds and completed 12 L-BFGS
iterations (30 value/gradient evaluations) in about 0.021 seconds after
compilation.  Writing a small checkpoint at every iteration did not dominate
this run; the final checkpoint was 237 bytes.  These figures demonstrate that
the operational wrapper has no visible cell-wise Python loop, but they remain
machine-dependent and do not replace the full-network admission benchmark.

